\documentclass[12pt]{article}
\usepackage[margin=1in]{geometry}
\usepackage{enumitem}
\usepackage{titlesec}
\usepackage{parskip}
\usepackage{amsmath}
\usepackage{hyperref}
\usepackage{fontenc}

\title{\textbf{Olson 1965 Claude Guide}\\
\large Olson, Mancur. \textit{The Logic of Collective Action: Public Goods and the Theory of Groups}.\\
Boston: Harvard University Press, 1965.}
\author{}
\date{}

\begin{document}

\maketitle
\tableofcontents
\newpage

%--------------------------------------------------
\section{Central Claims}
%--------------------------------------------------

\begin{itemize}[leftmargin=*, itemsep=4pt]
    \item The book's core puzzle: rational, self-interested individuals who share a common interest in some collective good will \textbf{not} voluntarily act to achieve that common interest, even when they agree the group would be better off if they did. This directly contradicts the ``naive'' assumption, implicit in pluralist group theory (Bentley, Truman), that shared interest is sufficient to generate group action.
    \item The mechanism is the \textbf{free-rider problem}: because the benefits of group action are typically a \textbf{public good}---\textbf{non-excludable} (members cannot be denied the benefit even if they didn't contribute) and characterized by \textbf{jointness of supply} (one member's consumption does not diminish another's)---any rational individual prefers to let others bear the costs of provision while still enjoying the benefit.
    \item \textbf{Group size is the decisive variable}. In small (``privileged'') groups, individual members may capture enough of the benefit that unilateral provision is rational, or the group is small enough that mutual monitoring and negotiation can sustain cooperation. In large (``latent'') groups, an individual's share of the collective benefit shrinks toward zero, monitoring becomes impossible, and the good will not be provided by voluntary, rational action alone.
    \item Large latent groups can only overcome the free-rider problem through \textbf{selective incentives}---positive or negative incentives applied specifically to individuals depending on whether they contribute to the collective effort, as distinct from the collective good itself (which cannot be selectively withheld).
    \item The \textbf{by-product theory}: mass-membership organizations that successfully mobilize large latent groups (above all labor unions and lobbies) typically do so not by providing the public good directly as their draw, but by ``bundling'' it with selectively-provided private goods (wages bargained via compulsory union shops, professional services, insurance, social solidarity)---the collective political good is a \emph{by-product} of an organization built to sell private benefits.
    \item This logic devastates both \textbf{orthodox Marxist class theory} (which assumes the proletariat will act as a class simply because it shares a class interest) and \textbf{simple pluralist/group theories of politics} (Bentley, Truman, and the ``group theory'' tradition), both of which conflate shared interest with organized action.
    \item Implication for interest group and lobbying theory: the universe of organized political interests is systematically \textbf{biased toward small, concentrated groups} (e.g., producers, industries, professions) and away from large, diffuse groups (e.g., consumers, taxpayers, the general public)---not because the latter's interests are less real or less shared, but because they face a harder collective action problem and typically lack pre-existing organizational infrastructure through which selective incentives can be distributed.
\end{itemize}

%--------------------------------------------------
\section{Core Concepts and Theoretical Apparatus}
%--------------------------------------------------

\begin{itemize}[leftmargin=*, itemsep=4pt]
    \item \textbf{Public/collective good}: any good such that if any person in a group consumes it, it cannot feasibly be withheld from other members of that group. Olson's definition is group-relative---a good can be ``public'' with respect to one group and irrelevant/private with respect to others.
    \item \textbf{Non-excludability}: the inability to exclude non-contributors from enjoying the good once it is provided; the structural source of the free-rider incentive.
    \item \textbf{Jointness of supply (non-rivalry)}: one person's consumption of the good does not reduce its availability to others.
    \item \textbf{Privileged groups}: groups small enough (or containing a member with a large enough stake) that at least one member finds it individually rational to provide the collective good unilaterally, even bearing the whole cost alone. Collective action is not automatic even here, but it is far more likely.
    \item \textbf{Intermediate groups}: groups where no single member has an incentive to provide the good alone, but the group is small enough that members can plausibly perceive their individual impact, negotiate, and monitor one another---so voluntary organization remains possible without selective incentives, though not guaranteed.
    \item \textbf{Latent groups}: large groups in which no individual's contribution has a noticeable effect on the group as a whole or on the burden/benefit of any other individual; individuals cannot be expected to act to further a common interest through voluntary, unorganized action, because they correctly perceive that their own contribution is both individually costly and collectively negligible.
    \item \textbf{The ``exploitation of the great by the small''}: in privileged and intermediate groups, the members with the largest stake in the collective good (the ``great'') will tend to bear a disproportionate share of the cost, since they have the greatest incentive to ensure provision, while smaller members free-ride on their efforts---an inversion of the intuition that large actors exploit small ones.
    \item \textbf{Selective incentives}: incentives that can be applied to individuals selectively (rewarding contributors, penalizing non-contributors), as opposed to the collective good itself which cannot be so targeted. May be \textbf{positive} (social status, camaraderie, private services, professional certification, insurance) or \textbf{negative} (social ostracism, compulsory dues, the union ``closed shop'' or ``union shop'').
    \item \textbf{The by-product theory of large pressure groups}: large-scale political organizations (unions, professional associations, farm lobbies) achieve mass membership and thereby political capacity not primarily through appeals to the collective political good but through provision of selectively-distributed private goods; the organization's lobbying/political activity is a by-product of an organizational apparatus built and sustained for other reasons.
    \item \textbf{Critique of the ``naive'' group theory}: the assumption, common to Bentley's and Truman's pluralist tradition, that individuals with common interests will spontaneously organize to further those interests. Olson shows this assumption is logically unsound once one takes rational self-interest and group size seriously.
    \item \textbf{Federal/coalition structure of large organizations}: large latent groups are often mobilized not directly but through federations of smaller, already-organized units (e.g., national labor federations built out of local unions, or agricultural lobbies built out of regionally organized cooperatives)---each sub-unit solves its own smaller collective action problem, and the aggregation solves the larger one.
\end{itemize}

%--------------------------------------------------
\section{Core Works Cited}
%--------------------------------------------------

\begin{itemize}[leftmargin=*, itemsep=4pt]
    \item \textbf{Arthur Bentley}, \textit{The Process of Government} (1908)---foundational group theory of politics; Olson's central foil for the ``naive'' assumption that shared group interest automatically yields group action.
    \item \textbf{David Truman}, \textit{The Governmental Process} (1951)---the mid-century pluralist synthesis of group theory that Olson targets most directly; Truman's account of interest-group formation as a natural response to shared interest is precisely what the logic of collective action refutes.
    \item \textbf{Karl Marx}, \textit{Das Kapital} and related writings on class consciousness and class action---Olson's critique of Marxist theory targets the assumption that a shared objective class interest (the proletariat's interest in revolution) will translate into class-conscious collective action; Olson argues Marx never solves the free-rider problem within the working class.
    \item \textbf{Paul Samuelson}, ``The Pure Theory of Public Expenditure'' (1954, \textit{Review of Economics and Statistics})---the canonical modern formalization of the public goods concept (non-rivalry, non-excludability) that Olson adapts to political and organizational contexts.
    \item \textbf{Kenneth Arrow}, \textit{Social Choice and Individual Values} (1951)---part of the broader rational-choice/welfare-economics toolkit Olson draws on; both share a commitment to formal, individual-level rationality as the basis for social theory.
    \item \textbf{Anthony Downs}, \textit{An Economic Theory of Democracy} (1957)---a close contemporary in applying economic/rational-choice reasoning to political behavior (voting, party competition); Olson extends the same logic to group formation and lobbying.
    \item \textbf{Selig Perlman}, \textit{A Theory of the Labor Movement} (1928), and other labor-economics historians---sources for Olson's extended case study of American labor union growth and the historical role of compulsory membership (the closed/union shop) in solving unions' collective action problem.
    \item \textbf{John R. Commons} and the ``Wisconsin School'' of labor economics---background institutional history of American unionism that Olson draws on and reinterprets through the free-rider/selective-incentive lens.
\end{itemize}

%--------------------------------------------------
\section{Chapter 1: A Theory of Groups and Organizations}
%--------------------------------------------------

\begin{itemize}[leftmargin=*, itemsep=4pt]
    \item Olson opens by stating the puzzle plainly: it is \emph{not} true, as ``the traditional theory of groups'' or pressure-group pluralism assumes, that groups of individuals with common interests usually attempt to further those interests. Common interest is necessary but not sufficient for group action.
    \item Distinguishes the ``furtherance of the common interests of a group'' from purely individual, uncoordinated behavior---the entire book concerns whether and when organizations or collective mechanisms emerge to pursue goods that benefit the group \emph{as a whole}.
    \item Introduces the formal logic: if a collective good is provided, it benefits every member of the group whether or not that member helped pay for it (non-excludability). A rational, self-interested individual who is not subject to coercion or a separate incentive will not contribute to the group's welfare, because his own effort has a negligible or unnoticeable effect on the outcome and full enjoyment of the benefit does not require his contribution.
    \item Uses simple economic logic (marginal cost vs.\ marginal share of collective benefit) to show that the ``rational'' amount of individual contribution to a large group's collective good is often zero, even for members who value the good highly and would prefer, in the aggregate, that the group succeed.
    \item Explicitly separates his argument from questions of altruism, ideology, or ``civic-mindedness''---he is not arguing individuals never contribute (some do, for non-narrowly-economic reasons), but that pure economic self-interest alone predicts systematic under-provision, and this baseline matters enormously for institutional design and lobbying behavior.
    \item Sets up the book's organizing variable: \textbf{group size}, which determines (a) whether any single member's share of the benefit is large enough to justify unilateral provision, and (b) whether the group is small enough for members to coordinate, bargain, and monitor each other directly.
\end{itemize}

%--------------------------------------------------
\section{Chapter 2: Group Size and Group Behavior}
%--------------------------------------------------

\begin{itemize}[leftmargin=*, itemsep=4pt]
    \item Formalizes the size typology: \textbf{privileged groups} (at least one member has enough at stake to provide the good alone), \textbf{intermediate groups} (no member can unilaterally justify provision, but the group is small enough that each member's contribution is perceptible and negotiation/monitoring are feasible), and \textbf{latent groups} (so large that individual contributions are imperceptible and organization requires external inducement).
    \item Derives the ``exploitation of the great by the small'': within privileged and intermediate groups, members with disproportionately large stakes in the collective good will tend to supply a disproportionate share of it, since their personal benefit alone may justify the cost, while smaller members can free-ride on the large member's provision. Applies this to international relations (e.g., large powers bearing disproportionate shares of alliance/defense burdens relative to small allies---a point later taken up extensively in the public-goods literature on NATO burden-sharing).
    \item Shows that as group size increases (holding total group benefit roughly constant), the likelihood of voluntary, unorganized collective action decreases, because (a) each member's share of the benefit shrinks, (b) the cost of organizing/negotiating among members rises, and (c) the ability to detect and sanction free-riders falls.
    \item Introduces \textbf{social pressure} as a potential selective incentive in intermediate groups (where members interact face-to-face and can observe each other's contributions), but argues this mechanism breaks down as group size grows and anonymity increases---social pressure cannot substitute for organization in truly large, latent groups.
    \item Distinguishes the costs of organizing a group (bargaining, information, monitoring) as themselves a kind of collective good subject to the same free-rider logic---which is part of why latent groups so often remain unorganized even when an entrepreneur or potential leader could, in principle, benefit all members by organizing them.
\end{itemize}

%--------------------------------------------------
\section{Chapter 3: The Traditional Theory of Labor Unions and Union Growth}
%--------------------------------------------------

\begin{itemize}[leftmargin=*, itemsep=4pt]
    \item Uses the American labor movement as the book's extended empirical case study of large-group collective action. Wages and working conditions won through collective bargaining are public goods with respect to the workers in a bargaining unit: a wage increase benefits union and non-union workers alike, and cannot be withheld from non-members once negotiated.
    \item Critiques ``traditional'' labor economics accounts of union growth (institutional/historical narratives, appeals to worker solidarity or class consciousness) for failing to explain \emph{why} workers would individually choose to pay dues and bear the risks of union membership (strikes, employer retaliation) rather than free-ride on a union's bargained gains.
    \item Documents empirically that union membership and growth historically track not moments of maximal worker solidarity but the availability and enforcement of \textbf{compulsory membership arrangements}---the closed shop and union shop---which convert the collective good of union representation into something workers cannot access without paying dues (a negative selective incentive: exclusion from employment, or at least from certain jobs, absent membership).
    \item Where compulsory membership is legally unavailable or weakly enforced (right-to-work contexts), Olson predicts and finds weaker, more fragile union organization, consistent with the free-rider logic rather than with variation in worker sentiment or solidarity.
    \item This case study is the empirical proving ground for the \textbf{by-product theory} developed fully in the following chapters: unions that survive and grow as large-scale organizations typically do so via coercive or selective mechanisms (compulsory dues, closed shops) rather than voluntary appeals to the shared interest in better wages.
\end{itemize}

%--------------------------------------------------
\section{Chapter 4: Orthodox Theories of the State, of Class, and of Pressure Groups}
%--------------------------------------------------

\begin{itemize}[leftmargin=*, itemsep=4pt]
    \item Extends the critique of naive group theory to two major bodies of political thought: \textbf{Marxist class theory} and \textbf{pluralist ``group theory'' of the state} (Bentley, Truman, and the broader American pluralist tradition).
    \item Against Marx: argues that ``class interest'' alone cannot explain working-class revolutionary action, because the benefits of a successful revolution (a transformed economic/political order) would be a public good for the entire class, including those who did not participate in or risk anything for the struggle. Marx's theory, on Olson's reading, silently assumes away the free-rider problem within the proletariat---Olson finds this a fundamental, previously unrecognized weakness in Marxist theory.
    \item Against Bentley/Truman-style pluralism: shows that treating ``interest groups'' as the natural, near-automatic organizational expression of shared social interests cannot be right, because it cannot explain the systematic \emph{absence} of organization among large, diffuse interests (consumers, unorganized labor, taxpayers, the poor) that are, by hypothesis, at least as ``interested'' in political outcomes as small, well-organized producer or professional groups.
    \item Notes the resulting asymmetry: the interest-group universe we actually observe is heavily populated by small, concentrated-interest groups (industries, professions, single firms) precisely because such groups are privileged or intermediate rather than latent, and/or because they can piggyback on pre-existing organizations (trade associations, professional bodies) that already solve smaller-scale collective action problems.
    \item This chapter is the pivot from ``why voluntary collective action in large groups is rare'' to ``what conditions allow large-group organizations to exist anyway''---setting up the by-product theory of Chapter 5.
\end{itemize}

%--------------------------------------------------
\section{Chapter 5: A Theory of the Political Process --- The By-Product Theory}
%--------------------------------------------------

\begin{itemize}[leftmargin=*, itemsep=4pt]
    \item Develops the \textbf{by-product theory} of large pressure groups: mass-membership political organizations representing large latent groups typically arise as a \emph{by-product} of an organization built primarily to sell members some other, selectively-distributed good---insurance, professional certification and services, trade publications, social fellowship, or (as in unions) compulsory workplace representation.
    \item The organization's lobbying and political activity is then financed out of resources (dues, staff, infrastructure) that were raised for these other, non-collective-action purposes---meaning the large group's political voice is essentially subsidized by an unrelated selective-incentive structure, not sustained directly by members' shared political interest.
    \item Explains the historical pattern whereby many large American lobbies (farm organizations, professional associations such as the American Medical Association and American Farm Bureau, veterans' organizations) grew out of organizations initially built to provide services, insurance, or professional/technical support, with political lobbying emerging afterward as a natural extension once the organization already had a mass, dues-paying membership.
    \item Implication: \textbf{political entrepreneurship} and pre-existing organizational infrastructure matter enormously for whether a large latent group's interests get represented at all; interests without such infrastructure (the classic case being consumers or the unorganized poor) remain latent and politically underrepresented even when their aggregate stake in policy outcomes is very large.
    \item This is Olson's answer to the pluralist assumption that the interest-group system reflects the universe of social interests roughly proportionally: it does not, and the by-product theory explains \emph{why} representation is systematically skewed toward interests that can be bundled with a selective-incentive-generating organization.
\end{itemize}

%--------------------------------------------------
\section{Chapter 6: The `Special Interest' Theory Reconsidered}
%--------------------------------------------------

\begin{itemize}[leftmargin=*, itemsep=4pt]
    \item Returns to labor unions and other ``special interest'' organizations to reconsider their political role in light of the by-product theory: unions' political lobbying arms (political action, legislative advocacy) are sustained by dues collected via compulsory-membership arrangements justified on non-political (bargaining/representation) grounds.
    \item Addresses the normative and institutional implications: because large-group political representation typically depends on coercive or selective mechanisms rather than the pure voluntary expression of shared interest, the strength of an interest in the political process is a poor guide to the number of people who share that interest or its aggregate importance.
    \item Considers government's role directly: legal support for compulsory unionism (e.g., protections for union/closed shops under U.S.\ labor law) is itself a state-provided selective-incentive structure that helps large latent groups (workers) overcome the free-rider problem that would otherwise prevent effective organization---an early and influential statement of how legal/institutional design can substitute for, or complement, private selective incentives.
    \item Closes the book's core argument by generalizing: the pattern of interest-group activity observed in any political system will reflect not the distribution of shared interests in the population but the distribution of \emph{organizational capacity}---itself a function of group size, the availability of selective incentives, and pre-existing organizational infrastructure.
\end{itemize}

%--------------------------------------------------
\section{Methodological Contributions and Limitations}
%--------------------------------------------------

\begin{itemize}[leftmargin=*, itemsep=4pt]
    \item \textbf{Foundational contribution to rational choice political science}: Olson is among the first to systematically apply microeconomic reasoning (individual rationality, marginal cost-benefit calculation) to the study of group formation and political organization, helping launch the rational-choice/public-choice research program in political science alongside Downs and Arrow.
    \item \textbf{Formal clarity, informal application}: the core logic is presented with simple, largely verbal/graphical economic reasoning rather than fully specified game-theoretic models; later scholars (Hardin, Chamberlin, and the broader public-choice/game-theory literature) formalize and refine Olson's arguments, including relaxing some of his simplifying assumptions (e.g., that individuals are purely and narrowly self-interested, or that group members are homogeneous in preferences and resources).
    \item \textbf{The group-size typology is somewhat under-specified at the margins}: the boundaries between ``privileged,'' ``intermediate,'' and ``latent'' groups are not given precise formal criteria, and later work (e.g., on step-goods, thresholds, and tipping models of collective action) refines when exactly organization becomes infeasible.
    \item \textbf{Empirical grounding is narrower than the theoretical claims}: the primary sustained case study is American labor unions; the broader claims about lobbying, class action, and interest-group systems are argued largely by logical demonstration and illustrative example rather than systematic comparative empirical testing across many large-group cases.
    \item \textbf{Understates non-material motivations}: critics (and later Olson-influenced scholars) note the theory brackets ideology, norms, and other-regarding preferences almost entirely; this yields sharp, falsifiable predictions but likely overstates the prevalence of pure free-riding relative to real-world collective action, which very often does occur without selective incentives (a tension later engaged directly by Ostrom 1990).
    \item \textbf{Static equilibrium framing}: the analysis is largely comparative-static (does an equilibrium of provision exist given group size and incentive structure?) rather than dynamic, and has less to say about how organizations are founded in the first place (the ``entrepreneurial'' moment) than about how they are sustained once they exist---an issue picked up by subsequent political entrepreneurship literature.
\end{itemize}

%--------------------------------------------------
\section{Connections to the Broader Literature}
%--------------------------------------------------

\begin{itemize}[leftmargin=*, itemsep=4pt]
    \item \textbf{Ostrom 1990 (\textit{Governing the Commons})}: The most direct and important complication of Olson's argument in the comps reading list. Ostrom marshals extensive empirical evidence that self-organized, voluntary collective action to manage common-pool resources \emph{does} regularly succeed in real-world latent/intermediate groups, without external coercion or purely material selective incentives---through communication, graduated sanctions, monitoring, and locally-crafted institutional rules. Ostrom does not reject Olson's logic so much as show that the conditions under which it binds (anonymity, absence of repeated interaction, absence of endogenous institution-building) are narrower than Olson's presentation implies, and that groups can often design their way out of the free-rider trap.
    \item \textbf{Putnam 1993 (\textit{Making Democracy Work})}: Putnam explicitly engages Olson (his guide cites Olson 1965 as ``the baseline problem that social capital solves''). Where Olson is pessimistic about voluntary large-group cooperation absent selective incentives, Putnam argues that accumulated \textbf{social capital}---norms of generalized reciprocity, dense horizontal networks, and social trust---can sustain cooperation and solve collective action problems without coercion or material side-payments, essentially proposing a historical/sociological substitute for Olson's selective incentives.
    \item \textbf{Huntington 1968 (\textit{Political Order in Changing Societies})}: Both scholars are concerned with the gap between diffuse social interests/demands and organized political capacity. Huntington's account of institutionalization as what channels and absorbs mobilized participation resonates with Olson's point that organizational infrastructure, not raw interest, determines political voice; Huntington's ``praetorian'' societies, in Olson's terms, might be read as societies where participation outpaces the organizational (selective-incentive) structures needed to channel it into stable interest representation.
    \item \textbf{Dahl 1972 (\textit{Polyarchy})}: Dahl's conditions for polyarchy include organizational pluralism, but Dahl (in the Bentley/Truman tradition Olson critiques) is less attentive to \emph{why} some interests organize and others do not. Olson supplies the missing micro-foundation: polyarchic competition among organized groups will systematically underrepresent large diffuse interests relative to concentrated ones, a point with real implications for how ``pluralist'' Dahl's polyarchies actually are in practice.
    \item \textbf{Moore 1966 (\textit{Social Origins of Dictatorship and Democracy})}: Moore's class-coalition analysis (landed elites, peasants, bourgeoisie) implicitly assumes classes can act collectively in pursuit of shared class interests during critical historical junctures. Olson's critique of Marxist class-action theory poses a direct challenge to this style of argument: Moore's analysis would benefit from, but generally does not supply, an account of \emph{how} large class actors overcome free-rider problems to act as coherent political agents (e.g., why peasants revolt in some cases and not others might hinge on local organizational/selective-incentive structures rather than shared objective class interest alone).
    \item \textbf{Cox 1997 (\textit{Making Votes Count})}: Cox's analysis of strategic coordination among voters and elites under different electoral rules is a close cousin of Olson's logic: voters and parties face their own collective action/coordination problems (e.g., strategic voting to avoid vote-splitting), and electoral institutions function partly as selective-incentive-like structures (rewarding coordinated behavior, penalizing miscoordination via wasted votes) that determine whether large groups of like-minded voters can act cohesively.
    \item \textbf{Lijphart 1999 (\textit{Patterns of Democracy})}: Lijphart's consensus democracies rely on encompassing, well-organized peak associations (e.g., corporatist labor and business federations) bargaining directly with government. Olson's by-product theory helps explain how such peak associations achieve and sustain mass membership in the first place (often via compulsory or near-compulsory membership structures), which is a precondition for the kind of encompassing interest intermediation Lijphart's consensus model depends on.
    \item \textbf{Weber 1930 (\textit{The Protestant Ethic and the Spirit of Capitalism})}: Where Weber roots economic behavior in internalized ethical/religious motivation (a non-material, norm-based account of action), Olson's framework is its rational-choice mirror image, bracketing normative/ideational motivation almost entirely in favor of narrow material self-interest. The contrast is a useful methodological pairing for comps: Weberian verstehen-based causal narrative vs.\ Olsonian deductive rational-choice modeling as two poles of explanatory strategy in the field.
    \item \textbf{Huber and Shipan 2002 (\textit{Deliberate Discretion?})}: Huber and Shipan's analysis of legislative delegation to bureaucracies concerns how principals (legislatures) solve their own collective and informational problems in writing statutes; Olson's account of organized interests (who lobby legislatures and agencies) supplies part of the demand-side political environment in which discretion and statutory control choices are made---concentrated, well-organized interests are more likely to secure detailed, favorable delegation language or effective oversight capacity.
    \item \textbf{Powell 2000 (\textit{Elections as Instruments of Democracy})}: Powell's concern with how electoral rules translate citizen preferences into government policy assumes citizens and groups can express and aggregate preferences; Olson's logic reminds us that the preference \emph{expression} stage is itself filtered by organizational capacity, so majoritarian vs.\ proportional systems may differentially advantage organized minorities (Olson's privileged/intermediate groups) versus diffuse majorities (latent groups) in ways Powell's aggregation-focused framework does not directly address.
    \item \textbf{Stokes 2013 (\textit{Brokers, Voters, and Clientelism})}: Clientelist machine politics can be read through an Olsonian lens as a selective-incentive solution to a mobilization problem: parties cannot rely on diffuse programmatic appeals (a public good among supporters) to guarantee turnout/support, so they use brokers to distribute targeted, excludable private benefits (a negative/positive selective incentive analogous to Olson's union shop) to solve their own collective action problem of mobilizing large, latent electorates.
    \item \textbf{Svolik 2012 (\textit{The Politics of Authoritarian Regimes})}: Svolik's analysis of authoritarian power-sharing and the ruling coalition's collective action problems (how to prevent the dictator from being overthrown, and how the dictator prevents elite collusion against him) is a direct descendant of Olson's framework applied to small, privileged groups (ruling coalitions) where monitoring and side-payments (Olson's selective incentives) are feasible, in contrast to the mass public, which typically remains a latent group unable to coordinate a revolutionary challenge absent extraordinary circumstances.
    \item \textbf{Carey 2009 (\textit{Legislative Voting and Accountability})}: Carey's analysis of party discipline and legislative voting cohesion concerns how party leaders solve an internal collective action problem---inducing individual legislators (who might individually prefer to defect for local benefit) to support the party line---via selective incentives (committee assignments, campaign support, leadership positions) that are recognizably Olsonian in structure, applied to the legislative party as the ``group.''
\end{itemize}

\end{document}
